% [原文 p.211]
\par\medskip\noindent\textbf{（定理 5.35 的证明，续）}\quad 若 $k\ge3$，则 $Y=B_0X$ 满足条件。

其次，假定 $k=1$ 或 $k=2$。此时 $H$ 的嵌入维数为 $3$。因此 $X$ 的嵌入维数为 $4$。把 $X$ 换为 $0$ 的小解析近旁，即可把 $X$ 看作由方程 $f(x,y,z,t)=0$ 定义的 $\C^4$ 的超曲面。

使用定理 4.57 所给 $H$ 的方程。该方程表明 $X$ 有二重点。因此可把方程写成 $x^2+g(y,z,t)=0$。若 $\operatorname{mult}_0g=2$，则方程还可变形为 $x^2+y^2+h(z,t)$。其超平面截面 $t=\lambda z$ 成为 $A$ 型 Du Val 奇点，与假定矛盾。因此 $\operatorname{mult}_0g\ge3$。

若 $\operatorname{mult}_0g\ge4$，则引入权 $\operatorname{wt}(x,y,z,t)=(2,1,1,1)$，并令 $f\colon Y\to X$ 为相应的加权爆破。

若 $\operatorname{mult}_0g=3$，则如 4.25 所述，方程可表示为 $x^2+y^3+yq_4(z,t)+q_6(z,t)$。若 $\operatorname{mult}_0q_4<4$ 或 $\operatorname{mult}_0q_6<6$，则由 4.25 中的第 6 步，超平面截面 $t=\lambda z$ 为 Du Val 奇点。于是 $\operatorname{mult}_0q_4\ge4$ 且 $\operatorname{mult}_0q_6\ge6$。定义权 $\operatorname{wt}(x,y,z,t)=(3,2,1,1)$，并令 $f\colon Y\to X$ 为相应的加权爆破。

无论哪种情形，假定 (1)(2) 都容易验证。由定义，$Y=\operatorname{Proj}_X\bigoplus_n\mathfrak{m}^w(n)$（理想 $\mathfrak{m}^w(n)$ 在定义 4.56 中定义）。因此 (3) 由下述结果得到。

\par\medskip\noindent\textbf{主张 5.37}\quad 设 $f(x_1,x_2,x_3,x_4)\in\C[[x_1,x_2,x_3,x_4]]$ 是幂级数。假定下列之一成立。
\begin{enumerate}[label=(\arabic*),leftmargin=*]
\item $f=(\text{单位})\cdot\bigl(x_1^2+g(x_2,x_3,x_4)\bigr)$，且 $\operatorname{mult}_0g\ge4$。此时，权由 $\operatorname{wt}(x_1,x_2,x_3,x_4)=(2,1,1,1)$ 给定。
\item $f=(\text{单位})\cdot\bigl(x_1^2+x_2^3+x_2q_4(x_3,x_4)+q_6(x_3,x_4)\bigr)$，且 $\operatorname{mult}_0q_r\ge r$。此时，权由 $\operatorname{wt}(x_1,x_2,x_3,x_4)=(3,2,1,1)$ 给定。
\end{enumerate}
设 $\varphi$ 是 $\C[[x_1,x_2,x_3,x_4]]$ 的 $\C$-自同构，且 $\varphi(f)$ 也具有同样形式的表示。（也就是说，(1) 的情形下假定 $\varphi(f)=(\text{单位})\cdot\bigl(x_1^2+g'(x_2,x_3,x_4)\bigr)$ 且 $\operatorname{mult}_0g'\ge4$。(2) 的情形也同样。）
% [原文 p.212]
此时，$\varphi\bigl(\mathfrak{m}^w(n)\bigr)=\mathfrak{m}^w(n)$（$\forall n\ge0$）。特别地，$\mathfrak{m}^w$ 只由 $f$ 决定，不依赖坐标的选取。

\begin{proof}
记 $\varphi_i:=\varphi(x_i)$。若 $\sum a_iw(x_i)\ge n$，则须证 $w\bigl(\prod_i\varphi_i^{a_i}\bigr)\ge n$。也就是说，$w(\varphi_i)\ge w(x_i)$ 是应当证明的主张。

关于 (1)，唯一须证明之处是 $\varphi_1=(\text{常数})x_1+(\text{次数}\ge2\text{ 的项之和})$。设它不成立。也就是说，例如假定 $\varphi_1$ 中 $x_2$ 的系数为 $c$（$\ne0$）。此时 $\varphi_1^2$ 中 $x_2^2$ 的系数为 $c^2$。又因 $\operatorname{mult}g\ge4$，$\varphi_1^2+g(\varphi_2,\varphi_3,\varphi_4)$ 中 $x_2^2$ 的系数仍为 $c^2$，这与 $\varphi(f)=(\text{单位})\cdot\bigl(x_1^2+g'(x_2,x_3,x_4)\bigr)$（$\operatorname{mult}_0g'\ge4$）矛盾。用同样的论证也可解决 (2)，但这留给读者。
\end{proof}

% [Tutor Note: 主张 5.37 的要旨是：权 $\mathrm{wt}$ 一旦由 $f$ 的主项确定，就与坐标的选取无关。证明只看最低次项，故只需检查 $\varphi_1$ 的一次项；高次项对 $x_2^2$ 的系数没有贡献。]

下面把这些结果的若干推论作为系叙述。

\begin{corollary}{5.38}
设 $0\in X$ 是三维正规奇点。此时，$0\in X$ 终端且指数为 $1$，当且仅当 $0\in X$ 是孤立 cDV 奇点。
\end{corollary}

关于指数大于 $1$ 的终端奇点，把推论 5.38 与推论 5.21 结合起来，立即得到下述结果。

\begin{corollary}{5.39}
若 $0\in X$ 是三维终端奇点，则它是孤立 cDV 奇点（或非奇异点）$0\in Y$ 由在原点之外自由作用的循环群作商。
\end{corollary}

\begin{corollaryp}{5.40}
设 $X$ 是三维簇，且只有至多标准奇点。此时，$X$ 中非 cDV 的奇点只有有限个。
\end{corollaryp}

\begin{proof}
由推论 5.18，$K_X$ 非 Cartier 的点只有有限个。若 $K_X$ 在 $x$ 处为 Cartier，而 $x\in X$ 非 cDV，则由定理 5.35，存在位于 $X$ 上方的例外除子 $E$，满足 $\operatorname{center}_XE=\{x\}$ 且 $a(E,X)=0$。由命题 2.36，位于 $X$ 上方且满足 $a(E,X)=0$ 的除子 $E$ 至多有限个。因此非 cDV 的点至多有限个。
\end{proof}

\begin{corollaryp}{5.41}
设 $0\in X'$ 是三维标准奇点，其指数为 $r$。设 $p\colon X\to X'$ 是指数 $1$ 覆盖。设 $X$ 没有 cDV 奇点。此时，存在 $\mathfrak{m}_{0,X'}$-准素
% [原文 p.213]
理想 $I\subset\cO_{X'}$，具有下述性质。$Y':=B_IX'$ 正规，$g\colon Y'\to X'$ 双有理，例外除子 $E\subset Y'$ 不为 $0$，且 $K_{Y'}\equiv p^*K_{X'}$ 成立。
\end{corollaryp}

\begin{proof}
设 $f\colon Y\to X$ 是定理 5.35 所导出的加权爆破。设 $N\in\Z_{>0}$ 是各坐标变量的权的公倍数，则存在唯一的有效 $f$-例外除子 $E\subset Y$，使 $E\sim\cO_Y(-N)$。设 $G$ 是 $X/X'$ 的伽罗瓦群。由定理 5.35，理想 $\mathfrak{m}^w(n)$ 是 $G$-不变的。于是得到 $Y$ 上的 $G$-作用，且 $E$ 是 $G$-不变的。

由于尚未证明 $Y$ 正规，设 $\bar f\colon\bar Y\to Y$ 是 $Y$ 的正规化，$\bar E$ 是 $E$ 在 $\bar Y$ 中的拉回。此时 $-\bar E$ 为 $\bar f$-丰富，故可取 $s>0$，使 $J:=\bar f_*\cO_{\bar Y}(-s\bar E)$ 生成 $\bigoplus_n\bar f_*\cO_{\bar Y}(-ns\bar E)$。把 $J$ 自然地看作 $\cO_X$ 的 $\mathfrak{m}_{0,X}$-准素理想，便得到 $Y=B_JX$。

记 $I:=J^G\subset\cO_{X'}$。此时 $B_IX'$ 是 $Y$ 在 $G$ 作用下的商，$I$ 是 $\mathfrak{m}_{0,X'}$-准素理想。由于 $Y$ 是正规簇在群作用下的商，$Y'$ 正规。

唯一剩下的是证明 $K_{Y'}\equiv p^*K_{X'}$。无论哪种情形，都可用例外除子 $E'_i\subset Y'$ 表示为 $K_{Y'}\equiv p^*K_{X'}+\sum a(E'_i,X')E'_i$。设 $E_i\subset Y$ 是支配 $E'_i$ 的 $f$-例外除子。命题 5.20 中已给出 $a(E_i,X)+1=m\bigl(a(E'_i,X')+1\bigr)$，其中 $m$ 是 $E_i$ 处的分歧指数。由于 $X'$ 为标准奇点，$a(E'_i,X')\ge0$；又由定理 5.35，$a(E_i,X)=0$。由此得到 $a(E'_i,X')=0$。
\end{proof}

虽然偏离正题，但请注意：定理 5.34 的证明也可用作 [Elk78] 中下述定理的证明。

\begin{theoremp}{5.42}
设 $0\in X$ 是奇点，$0\in H\subset X$ 是 Cartier 除子。若 $0\in H$ 具有有理奇点，则 $0\in X$ 也具有有理奇点。
\end{theoremp}

\begin{proof}[证明概要]
因 $H$ 为 CM，由定理 5.10，$X$ 也为 CM。对图 (5.3) 所示情形，已证明 $\beta$ 与 $\gamma$ 为同构。因此 $\delta\colon f_*\omega_Y\to\omega_X$ 也为同构。由定理 5.10，$X$ 具有有理奇点。
\end{proof}

本节以下叙述三维终端奇点的详细分类结果。不过，
% [原文 p.214]
本书不需要这些结果。

指数 $1$ 覆盖非奇异的情形，即三维终端商奇点，已由 [MS84]、[Dan82] 用环面嵌入理论分类，完成者是 M. Frumkin。其核心的“终端引理”（terminal lemma）被查明是有关格点的 G.K. White 定理这一古典结果。关于这些，[石井97, §7.3]、[小田85, §1.7] 中有解说。

指数 $1$ 覆盖具有奇点的情形，其按具体方程的分类由 [Mor85] 完成。充分性条件在 [KSB88] 中给出，其证明由 [Ste88] 简化。与 Du Val 奇点的覆盖的关系由 [Rei87b] 发现。详见 [Mor85] 与 [Rei87b]。

\begin{theorem*}{5.43}
设 $(0\in X)$ 是三维正规孤立奇点。设 $K_X$ 为 $\Q$-Cartier，指数为 $r$，$\pi\colon(\tilde0\in\tilde X)\to(0\in X)$ 是指数 $1$ 覆盖。此时，$1$ 的 $r$ 次根构成的群 $\mu_r$ 作用在 $\tilde X$ 上。
\begin{enumerate}[label=(\arabic*),leftmargin=*]
\item $(0\in X)$ 终端当且仅当含 $0$ 的一般元 $H\in|-K_X|$ 为 Du Val。
\item 表 5.1 是 $\tilde H:=\pi^*(H)$、$H$ 以及 $\mu_r$ 在 $\C^4$ 上作用的正合列表。
\end{enumerate}
\end{theorem*}
\begin{table}[htbp]
\centering
\tabtitle{5.1}{三维终端奇点与指数 $1$ 覆盖的正合列表}
\begin{tabular}{cccc}
\toprule
名称 & $\tilde H\to H$ 的类型 & $r$ & 作用的类型 \\
\midrule
$\mathrm{cA}/r$ & $A_{k-1}\to A_{kr-1}$ & $r$ & $1/r(a,-a,1,0;0)$ \\
$\mathrm{cAx}/2$ & $A_{2k-1}\to D_{k+2}$ & $2$ & $1/2(0,1,1,1;0)$ \\
$\mathrm{cAx}/4$ & $A_{2k-2}\to D_{2k+1}$ & $4$ & $1/4(1,1,3,2;2)$ \\
$\mathrm{cD}/2$ & $D_{k+1}\to D_{2k}$ & $2$ & $1/2(1,0,1,1;0)$ \\
$\mathrm{cD}/3$ & $D_4\to E_6$ & $3$ & $1/3(0,2,1,1;0)$ \\
$\mathrm{cE}/2$ & $E_6\to E_7$ & $2$ & $1/2(1,0,1,1;0)$ \\
\bottomrule
\end{tabular}
\end{table}
表中，$1/r(a_1,\dots,a_4;b)$ 表示 $\mu_r$ 的生成元 $\xi$ 按 $(x_1,\dots,x_4;f)\mapsto(\xi^{a_1}x_1,\dots,\xi^{a_4}x_4;\xi^bf)$ 作用于坐标 $x_1,\dots,x_4$ 与 $\tilde X$ 的方程 $f$。

“$\tilde H\to H$”这一列，也与在奇点之外非分歧的 Du Val 奇点循环覆盖的正合列表一致。

由定理 5.28 注意：指数 $r$ 的终端奇点 $0\in X$ 的任意平坦形变，都由指数 $1$ 覆盖 $\tilde X$ 的平坦形变的 $\mu_r$ 商给出。超曲面
% [原文 p.215]
奇点的形变由定理 4.61 具体描述。从 [Mor85] 所给方程可以读出下述结果。

\begin{corollary*}{5.44}
设 $f\colon X\to S$ 是纤维维数为 $3$ 的平坦态射。取闭点 $s_0\in S$ 与 $0\in X_{s_0}$。设 $0\in X_{s_0}$ 为终端，指数为 $r$。此时，存在近旁 $0\in X^0\subset X$，使得对满足 $X^0_s\ne\varnothing$ 的所有 $s\in S$ 都有下列性质。
\begin{enumerate}[label=(\arabic*),leftmargin=*]
\item $X^0_s$ 只有终端奇点。
\item $X^0_s$ 的每个奇点的指数都是 $r$ 的约数，且 $X^0_s$ 至少含有一个指数 $r$ 的终端奇点。
\end{enumerate}
\end{corollary*}

% [Tutor Note: 定理 5.43 与表 5.1 是三维终端奇点理论的高峰：所有指数 $1$ 的三维终端奇点都归结为 cDV 型，而一般终端奇点只是它的循环商。表中“名称”一栏 cA、cAx、cD、cE 分别指示上同调型。]

\section{伴随反转}
\takeaway{本节证明比较 $(X,S+B)$ 与 $(S,B|_S)$ 差异系数的若干定理。核心是连通性定理 5.48 与伴随反转定理 5.50：边界上的 klt／lc 性等价于原对的 plt／lc 性。最后给出 dlt 对边界既约部分的性质。}

设 $(X,S+B)$ 是 $n$ 维对。其中 $S$ 是 Weil 除子，$B$ 是 $\Q$-除子。本节目的是证明若干结果，用它们比较 $(X,S+B)$ 与 $(S,B|_S)$ 的差异系数。这类定理在借维数归纳处理许多问题时不可或缺。本节最后把这些结果应用于 $(X,S+B)$ 为 dlt 的情形。

一如既往，全部在代数设定下叙述，但对解析空间也同样成立。

需要把定义 2.28 如下精细化。

\begin{definition}{5.45}
设 $(X,\Delta)$ 是对，$Z\subset S\subset X$ 是闭部分概形。引入下述定义。
\[
\begin{aligned}
&\operatorname{discrep}(\operatorname{center}\subset Z,X,\Delta)\\
&\quad:=\inf\Bigl\{a(E,X,\Delta)\ \Bigm|\ \substack{E\ \text{为例外的，且}\\ \operatorname{center}_XE\subset Z}\Bigr\},\\[4pt]
&\operatorname{discrep}(\operatorname{center}\cap S\subset Z,X,\Delta)\\
&\quad:=\inf\Bigl\{a(E,X,\Delta)\ \Bigm|\ \substack{E\ \text{为例外的，且}\\ \operatorname{center}_XE\cap S\subset Z}\Bigr\}.
\end{aligned}
\]
其中 $E$ 取遍位于 $X$ 上方的全部例外除子。若允许 $E$ 取遍 $X$ 上方的一切除子（含非例外者），则这些 totaldiscrep 版也同样
% [原文 p.216]
可以定义。
\end{definition}

\begin{propositionp}[{[K+92, 17.2]}]{5.46}
设 $X$ 是正规簇，$S$ 是余维 $2$ 且为 Cartier 的正规 Weil 除子，$Z\subset S$ 是闭子簇，$B=\sum b_iB_i$ 是 $\Q$-除子。假定 $K_X+S+B$ 为 $\Q$-Cartier。此时下列成立。
\[
\begin{aligned}
&\operatorname{totaldiscrep}(\operatorname{center}\subset Z,S,B|_S)\\
&\quad\ge\operatorname{discrep}(\operatorname{center}\subset Z,X,S+B)\\
&\quad\ge\operatorname{discrep}(\operatorname{center}\cap S\subset Z,X,S+B).
\end{aligned}
\]
\end{propositionp}

\begin{proof}
设 $f\colon Y\to X$ 是 $(X,S+B)$ 的对数奇点消解，$S':=f^{-1}_*S$。再作爆破，可使 $\Supp f^{-1}(S+B)$ 非奇异，且与 $S'$ 相交的 $f$-例外除子 $E_i$ 全都满足 $\operatorname{center}_XE_i\subset S$。表示为 $K_Y+S'\equiv f^*(K_X+S+B)+\sum e_iE_i$。由伴随公式得到
\[
K_{S'}=(K_Y+S')|_{S'},\qquad (K_X+S+B)|_S=K_S+B|_S,
\]
由此导出下式。
\[
K_{S'}\equiv f^*(K_S+B|_S)+\textstyle\sum e_i(E_i\cap S').
\]
由于 $S'$ 与 $f^{-1}_*B$ 不相交，若 $E_i\cap S'\ne\varnothing$，则 $E_i$ 为 $f$-例外的且 $\operatorname{center}_XE_i\subset S$。这表明：出现在态射 $S'\to S$ 中的差异系数，必等于 $X$ 上方中心含于 $S$ 的 $Y\to X$ 的例外除子的差异系数。（$E_i$ 虽为 $f$-例外的，但 $E_i\cap S'$ 未必为 $f|_{S'}$-例外的。这正是左端出现 totaldiscrep 的理由。）
\end{proof}

一般地，$f\colon Y\to X$ 的例外除子 $E_j$ 中，与 $S'$ 不相交者是存在的。因此，这类除子的差异系数与出现在 $S'\to S$ 中的差异系数之间没有明显的关系。

尽管如此，[Sho92, 3.3] 与 [K+92, 17.3] 猜想命题 5.46 中成立等号。

这一猜想（或同类结果与猜想）在假定关于 $X$ 的某事物以得到关于 $S$ 的结论时，以伴随（adjunction）之名被引用。反之，在假定关于 $S$ 的某事物以得到关于 $X$ 的结论时，称为伴随反转（inversion of adjunction）。

\begin{remark}{5.47}
本节反复使用的假定是“$S\subset X$ 为余维 $2$ 且 Cartier
% [原文 p.217]
的 Weil 除子”。在此假定下 $K_S=(K_X+S)|_S$ 成立。若 $S$ 不是余维 $2$ 的 Cartier 除子，则该等式需要修正项。加上修正项，余维 $2$ 且 Cartier 这一假定就不再必要。就应用而言，这一推广不可或缺，但本书不需要。关于这一点让给 [K+92]（特别是第 16、17 章）。
\end{remark}

% [Tutor Note: 伴随与伴随反转只差方向：前者由 $X$ 推 $S$，后者由 $S$ 推 $X$。命题 5.46 给出的是伴随方向的不等式；定理 5.50 才是真正的伴随反转，即把 $S$ 上的 klt 性反推为 $X$ 上沿 $S$ 的 plt 性。]

伴随反转的许多应用中重要的是 klt 或 lc 的情形。这些证明中，下面叙述的连通性定理起中心作用。

\begin{theoremp}[{连通性定理 [K+92, 17.4]}]{5.48}
设 $g\colon Y\to X$ 是固有双有理态射，$Y$ 非奇异，$X$ 正规。设 $D=\sum d_iD_i$ 是 $Y$ 上的 snc $\Q$-除子，$g_*D$ 有效，$-(K_Y+D)$ 为 $g$-nef，并如下定义 $A,F$。
\[
A=\sum_{i:d_i<1}d_iD_i,\qquad F=\sum_{i:d_i\ge1}d_iD_i.
\]
此时，$\Supp F$ 在 $g$ 的任意纤维的近旁连通。
\end{theoremp}

\begin{proof}
由定义，
\[
\lfloor-A\rfloor-\lfloor F\rfloor=K_Y-(K_Y+D)+\{A\}+\{F\}
\]
成立。因此由推论 2.68，
\[
R^1f_*\cO_Y\bigl(\lfloor-A\rfloor-\lfloor F\rfloor\bigr)=0.
\]
把 $g_*$ 作用于正合列
\[
0\to\cO_Y\bigl(\lfloor-A\rfloor-\lfloor F\rfloor\bigr)\to\cO_Y\bigl(\lfloor-A\rfloor\bigr)\to\cO_{\lfloor F\rfloor}\bigl(\lfloor-A\rfloor\bigr)\to0,
\]
得到下述满射。
\[
g_*\cO_Y\bigl(\lfloor-A\rfloor\bigr)\to g_*\cO_{\lfloor F\rfloor}\bigl(\lfloor-A\rfloor\bigr).
\]
设 $D_i$ 是 $A$ 的任意不可约分支。由于 $g_*D$ 有效，由此可知 $D_i$ 为 $g$-例外的，或 $d_i>0$。因此 $\lfloor-A\rfloor$ 是 $g$-例外的且有效，$g_*\cO_Y(\lfloor-A\rfloor)=\cO_X$ 成立。对某个 $x\in X$，假定 $\lfloor F\rfloor$ 在 $g^{-1}(x)$ 的近旁不连通，即 $\lfloor F\rfloor=F_1\sqcup F_2$。由此得到
\[
g_*\cO_{\lfloor F\rfloor}\bigl(\lfloor-A\rfloor\bigr)_{(x)}\cong g_*\cO_{F_1}\bigl(\lfloor-A\rfloor\bigr)_{(x)}\oplus g_*\cO_{F_2}\bigl(\lfloor-A\rfloor\bigr)_{(x)}.
\]
然而这些直和成分都不为 $0$。因此 $g_*\cO_{\lfloor F\rfloor}(\lfloor-A\rfloor)_{(x)}$ 不可能是 $\cO_{x,X}\cong g_*\cO_Y(\lfloor-A\rfloor)_{(x)}$ 的商。
\end{proof}

% [原文 p.218]
\begin{corollaryp}{5.49}
设 $g\colon Y\to X$ 是正规簇之间的固有双有理态射，$D=\sum d_iD_i$ 是 $Y$ 上的 $\Q$-除子，且 $g_*D$ 有效，$-(K_Y+D)$ 为 $g$-nef 的 $\Q$-Cartier。设 $Z\subset Y$ 是 $(Y,D)$ 非 klt 的点全体构成的部分集合。此时，$Z$ 在 $g$ 的任意纤维的近旁连通。
\end{corollaryp}

\begin{proof}
把定理 5.48 用于 $(Y,D)$ 的对数奇点消解即可。
\end{proof}

作为推论得到下述结果。三维情形由 [Sho92] 证明，一般维数情形由 [K+92, 17.6-7] 证明。

\begin{theoremp}[{伴随反转}]{5.50}
设 $X$ 正规，$S\subset X$ 是余维 $2$ 且 Cartier 的正规 Weil 除子。设 $B$ 是有效 $\Q$-除子，且 $K_X+S+B$ 为 $\Q$-Cartier。此时下列成立。
\begin{enumerate}[label=(\arabic*),leftmargin=*]
\item $(X,S+B)$ 在 $S$ 的近旁为 plt 当且仅当 $(S,B|_S)$ 为 klt。
\item 进一步假定 $B$ 为 $\Q$-Cartier 且 $S$ 为 klt。此时，$(X,S+B)$ 在 $S$ 的近旁为 lc 当且仅当 $(S,B|_S)$ 为 lc。
\end{enumerate}
\end{theoremp}

\begin{proof}
无论哪种情形，$\Rightarrow$ 都由命题 5.46 导出。

为证 (1)，设 $g\colon Y\to X$ 是 $(X,S+B)$ 的对数奇点消解，与记号 2.26 同样地记 $K_Y+D\equiv g^*(K_X+S+B)$。

如定理 5.48 定义 $A,F$，记 $S'=g^{-1}_*S$，$F=S'+F'$。取伴随得到 $K_{S'}=g^*(K_S+B|_S)-(A+F')|_{S'}$。$(X,S+B)$ 在 $S$ 的近旁为 plt 当且仅当 $F'\cap g^{-1}(S)=\varnothing$，这很明显。又，$(S,B|_S)$ 为 klt 当且仅当 $F'\cap S'=\varnothing$，也很明显。假定 $(S,B|_S)$ 为 klt。由定理 5.48，任意 $x\in S$ 都有开近旁 $U_x\subset X$，使 $(S'\cup F')\cap g^{-1}(U_x)$ 连通。因此 $F'\cap g^{-1}(U_x)=\varnothing$。让 $x\in S$ 变动，即得 (1)。

由命题 2.43，$(X,S+B)$ 为 lc 当且仅当对一切 $c<1$，$(X,S+cB)$ 为 plt。同样，$(S,B|_S)$ 为 lc 当且仅当对一切 $c<1$，$(S,cB|_S)$ 为 klt。于是由 (1) 导出 (2)。
\end{proof}

上述手法对理解 dlt 对的既约边界非常有用。

\begin{propositionp}{5.51}
设 $(X,\Delta)$ 是 dlt 对，$\Delta$ 有效。此时下列条件等价。
\begin{enumerate}[label=(\arabic*),leftmargin=*]
\item $(X,\Delta)$ 为 plt。
% [原文 p.219]
\item $\lfloor\Delta\rfloor$ 正规。
\item $\lfloor\Delta\rfloor$ 的各连通成分不可约。
\end{enumerate}
\end{propositionp}

\begin{proof}
假定 (1)。设 $g\colon Y\to X$ 是对数奇点消解，使 $\Delta':=g^{-1}_*\Delta$ 非奇异。
\[
K_Y+\Delta'\equiv g^*(K_X+\Delta)+E
\]
其中，由于 $(X,\Delta)$ 为 plt，$\lfloor E\rfloor$ 有效且为 $g$-例外的。考虑下述正合列。
\[
0\to\cO_Y\bigl(-\lfloor\Delta'\rfloor+\lfloor E\rfloor\bigr)\to\cO_Y\bigl(\lfloor E\rfloor\bigr)\to\cO_{\lfloor\Delta'\rfloor}\bigl(\lfloor E\rfloor|_{\lfloor\Delta'\rfloor}\bigr)\to0.
\]
由定义，$-\lfloor\Delta'\rfloor+\lfloor E\rfloor\equiv_g K_Y+\{\Delta'\}+(\lfloor E\rfloor-E)$，故由推论 2.68，$R^1f_*\cO_Y(-\lfloor\Delta'\rfloor+\lfloor E\rfloor)=0$ 成立。于是得到下述满射。
\[
g_*\cO_Y\bigl(\lfloor E\rfloor\bigr)\to g_*\cO_{\lfloor\Delta'\rfloor}\bigl(\lfloor E\rfloor|_{\lfloor\Delta'\rfloor}\bigr).
\]
由于 $\lfloor E\rfloor$ 有效且为 $g$-例外的，$\cO_X=g_*\cO_Y(\lfloor E\rfloor)$ 成立。设 $\lfloor\Delta\rfloor^n$ 是 $\lfloor\Delta\rfloor$ 的正规化，则下述映射
\[
g_*\cO_Y\bigl(\lfloor E\rfloor\bigr)=\cO_X\to\cO_{\lfloor\Delta\rfloor}\subset\cO_{\lfloor\Delta\rfloor^n}\subset g_*\cO_{\lfloor\Delta'\rfloor}\bigl(\lfloor E\rfloor|_{\lfloor\Delta'\rfloor}\bigr)
\]
的合成是满射。因此 $\cO_{\lfloor\Delta\rfloor}=\cO_{\lfloor\Delta\rfloor^n}$，特别地 $\lfloor\Delta\rfloor$ 正规。

(2)$\Rightarrow$(3) 显然。

假定 (3)。设 $Z\subset X$ 是使 $X\setminus Z$ 非奇异且 $\Delta|_{X\setminus Z}$ 为 snc 的最小闭部分集合，$E$ 是位于 $X$ 上方的任意例外除子。若 $\operatorname{center}_XE\subset Z$，则由定义 $a(E,X,\Delta)>-1$。若 $\operatorname{center}_XE\not\subset Z$，则由推论 2.31(3)，$a(E,X,\Delta)>-1$。
\end{proof}

\begin{corollaryp}{5.52}
设 $(X,\Delta)$ 是 dlt 对，$\Delta$ 有效。此时，$\lfloor\Delta\rfloor$ 的所有不可约成分都正规。
\end{corollaryp}

\begin{proof}
设 $S\subset\lfloor\Delta\rfloor$ 是不可约成分，$\Delta_1:=\Delta-S$。由命题 2.43，对一切有理数 $0<\epsilon\ll1$，存在与 $\Delta_1$ 无公共成分的有效 $\Q$-除子 $D$，使 $(X,\Delta-\epsilon\Delta_1+\epsilon D)$ 为 dlt。此时 $\lfloor\Delta-\epsilon\Delta_1+\epsilon D\rfloor=S$，由命题 5.51，$S$ 正规。
\end{proof}

% [Tutor Note: 推论 5.52 的技巧是：略去边界其余成分的 $\epsilon$ 份，同时加入与 $S$ 横截的扰动除子，使新 dlt 对的边界恰为 $S$，再套用命题 5.51。]

\begin{corollaryp}{5.53}
设 $(X,S+\Delta)$ 是 dlt 对，$\Delta$ 有效。设 $S$ 不可约、余维 $2$ 且 Cartier，并进一步假定 $\lfloor\Delta\rfloor=0$。此时，$\lfloor\Delta|_S\rfloor=0$。
\end{corollaryp}

\begin{proof}
$(X,S+\Delta)$ 为 plt，$S$ 正规（命题 5.51）。因此，
% [原文 p.220]
由定理 5.50，$(S,\Delta|_S)$ 为 klt。于是 $\lfloor\Delta|_S\rfloor=0$。
\end{proof}

\begin{remark}{5.54}
推论 5.53 对 lc 对 $(X,S+\Delta)$ 不成立。例如，$X=\C^2$，$S=(x=0)$，$\Delta=(1/2)(x+y=0)+(1/2)(x-y=0)$ 就是其反例。
\end{remark}

\begin{corollaryp}{5.55}
设 $(0\in X,\Delta)$ 是二维 dlt 对的芽，$\Delta$ 有效。此时下列之一成立。
\begin{enumerate}[label=(\arabic*),leftmargin=*]
\item $X$ 非奇异，且 $\Delta=\lfloor\Delta\rfloor$ 由在 $0$ 处横截相交的 $2$ 个不可约成分构成。
\item $(0\in X,\Delta)$ 为 plt，且 $\lfloor\Delta\rfloor$ 非奇异。
\end{enumerate}
\end{corollaryp}

\begin{proof}
假定条件 (1) 不成立。设 $f\colon Y\to X$ 是任意奇点消解，$E\subset Y$ 是满足 $f(E)=\{0\}$ 的例外除子。此时由假定 $a(E,X,\Delta)>-1$。因此 $(0\in X,\Delta)$ 为 plt，且由命题 5.51，$\lfloor\Delta\rfloor$ 正规，即非奇异。
\end{proof}

\begin{corollaryp}{5.56}
设 $(S,D+\Delta)$ 是 dlt 对，$D$ 为 Cartier 除子，$\Delta$ 为有效 $\Q$-除子。此时，在 $D$ 的任意余维 $1$ 的点处 $S$ 非奇异。特别地，$D$ 的各不可约成分在 $X$ 上余维 $2$ 且为 Cartier。
\end{corollaryp}

\begin{proof}
首先，通过局部化或取一般超曲面截面，只需处理 $X$ 为曲面的情形。由推论 5.55，$S$ 或 $D$ 非奇异。后一情形中，由于 $D$ 为 Cartier，$S$ 也非奇异。特别地，任意 Weil 除子都是 Cartier。
\end{proof}

\begin{corollary}{5.57}
设 $(0\in X,D+\Delta)$ 是 $0\in X$ 与 $0\in D$ 均非奇异、$\Delta$ 有效的二维对，且 $D$ 与 $\Delta$ 无公共成分。此时，$(0\in X,D+\Delta)$ 在 $0$ 的近旁为 plt（或 lc）当且仅当 $(D\cdot\Delta)_0<1$（或 $\le1$）。其中 $(D\cdot\Delta)_0$ 表示在 $0$ 处的局部交点数。
\end{corollary}

% [Tutor Note: 推论 5.57 把二维 lc 判据化成一个局部交点数，无须消解即可检验；命题 5.58 正是以此为基础的曲面版伴随反转。]

下面叙述的是 lc 曲面情形的伴随反转定理。

\begin{propositionp}{5.58}
设 $f\colon S\to T$ 是正规曲面芽 $(0\in T)$ 的奇点消解。设 $D\subset S$ 是 $f$-非例外的正规既约除子，$s\in D\cap f^{-1}(0)$，$\Delta$ 是 $S$ 上的 $\Q$-除子，且 $f_*\Delta$ 有效，$D\not\subset\Supp\Delta$。进一步假定下述 $2$ 个条件。
\begin{enumerate}[label=(\arabic*),leftmargin=*]
\item $K_S+D+\Delta\equiv0$。
\item $(D,\Delta|_D)$ 在 $s$ 处为 lc（这等价于 $(D\cdot\Delta)_s\le1$）。
\end{enumerate}
此时，$(S,D+\Delta)$ 为 lc。
\end{propositionp}
